% !TEX root = ../LosningsforslagTRES0312.tex
\chapter{Fluiddynamikk}
\label{LF:Chapter:fluiddynamikk}
Dette kapittelet inneholder løsningsforslag til oppgavene i Kapittel~\ref{komp-Chapter:fluiddynamikk}, \textit{Fluiddynamikk}, i kompendiet.

\oppgaver{Hvilket metall?}{k7oppg1}{%
%
En liten metallterning har sidekanter på \SI{3.0}{\centi\metre} og masse \SI{73}{\gram}.
\begin{enumerate}[a)]
    \item Finn volumet til terningen i \si{\centi\metre\cubed} og i \si{\metre\cubed}.
    \item Finn massetettheten til terningen. Bruk tabell~\komplabel{tab:massetetthet} til å finne ut hvilket metall terningen sannsynligvis er laget av.
    \item Hvor stor masse ville en like stor terning av gull hatt?
\end{enumerate}
%
%
}

\svar[title={Løsningsforslag}]{%
\noindent
\textbf{a)} Volumet til en terning er sidekanten opphøyd i tredje:
\[
    V = (\SI{3.0}{\centi\metre})^3 = \doubleunderline{\SI{27}{\centi\metre\cubed}}
\]
Siden $\SI{1}{\centi\metre} = \SI{1e-2}{\metre}$, er $\SI{1}{\centi\metre\cubed} = (\SI{1e-2}{\metre})^3 = \SI{1e-6}{\metre\cubed}$. Dermed er
\[
    V = 27\cdot\SI{1e-6}{\metre\cubed} = \doubleunderline{\SI{2.7e-5}{\metre\cubed}}
\]

\smallskip\noindent
\textbf{b)} Vi bruker definisjonen av massetetthet (se \textit{Massetetthet} i Kapittel~\ref{komp-Chapter:fluiddynamikk}):
\[
    \rho = \frac{m}{V} = \frac{\SI{73}{\gram}}{\SI{27}{\centi\metre\cubed}} = \SI{2.70}{\gram\per\centi\metre\cubed} \approx \doubleunderline{\SI{2.7e3}{\kilo\gram\per\metre\cubed}}
\]
I tabell~\komplabel{tab:massetetthet} ser vi at dette passer med \doubleunderline{aluminium}.

\smallskip\noindent
\textbf{c)} Gull har $\rho = \SI{19.3e3}{\kilo\gram\per\metre\cubed}$. En terning av gull med samme volum har massen
\[
    m = \rho\cdot V = \SI{19.3e3}{\kilo\gram\per\metre\cubed}\cdot\SI{2.7e-5}{\metre\cubed} = \SI{0.521}{\kilo\gram} \approx \doubleunderline{\SI{0.52}{\kilo\gram}}
\]
Gullterningen ville altså vært omtrent sju ganger så tung som aluminiumsterningen, selv om de er like store.
}

\oppgaver{Tegnestift}{k7oppg2}{%
%
Du trykker en tegnestift inn i en vegg med tommelen. Kraften fra tommelen på stiften er $F = \SI{20}{\newton}$. Hodet på tegnestiften har areal \SI{1.0}{\centi\metre\squared}, og spissen har areal \SI{0.10}{\milli\metre\squared}.
\begin{enumerate}[a)]
    \item Finn trykket mot tommelen.
    \item Stiften er i ro i det øyeblikket du trykker, så den presser på veggen med den samme kraften. Finn trykket fra spissen mot veggen.
    \item Forklar hvorfor stiften går inn i veggen, men ikke inn i tommelen.
\end{enumerate}
\textit{Hint: Husk at $\SI{1}{\centi\metre\squared} = \SI{1e-4}{\metre\squared}$ og $\SI{1}{\milli\metre\squared} = \SI{1e-6}{\metre\squared}$.}
%
%
}

\svar[title={Løsningsforslag}]{%
\noindent
Vi bruker definisjonen av trykk, $p = F/A$ (se \textit{Trykk} i Kapittel~\ref{komp-Chapter:fluiddynamikk}).

\smallskip\noindent
\textbf{a)} Arealet av hodet er $A = \SI{1.0}{\centi\metre\squared} = \SI{1.0e-4}{\metre\squared}$:
\[
    p_\text{tommel} = \frac{F}{A} = \frac{\SI{20}{\newton}}{\SI{1.0e-4}{\metre\squared}} = \doubleunderline{\SI{2.0e5}{\pascal}}
\]

\smallskip\noindent
\textbf{b)} Arealet av spissen er $A = \SI{0.10}{\milli\metre\squared} = 0.10\cdot\SI{1e-6}{\metre\squared} = \SI{1.0e-7}{\metre\squared}$:
\[
    p_\text{vegg} = \frac{F}{A} = \frac{\SI{20}{\newton}}{\SI{1.0e-7}{\metre\squared}} = \doubleunderline{\SI{2.0e8}{\pascal}}
\]

\smallskip\noindent
\textbf{c)} Kraften er like stor i begge ender av stiften, men spissen har et areal som er tusen ganger mindre enn hodet. Derfor er trykket mot veggen tusen ganger større enn trykket mot tommelen. Det store trykket er nok til at spissen trenger inn i veggen, mens trykket mot tommelen er så lite at det ikke gjør vondt.
}

\oppgaver{Vannkanne}{k7oppg3}{%
%
En vannkanne er fylt med vann som står i ro. Punkt A ligger inne i kanna, og punkt B ligger i tuten, i samme høyde som A (se figuren).

\begin{center}
\begin{tikzpicture}[scale=0.9, font=\small]
  % Vann (klippes ved vannflaten y = 2.4)
  \begin{scope}
    \clip (-1,-1) rectangle (5,2.4);
    \fill[oppgcol!18] (0,0) rectangle (2.4,3.0);
    \fill[oppgcol!18] (2.35,0.3) -- (3.75,3.0) -- (3.15,3.0) -- (2.35,1.5) -- cycle;
  \end{scope}
  \draw[oppgcol!70!black, thick] (0,2.4) -- (2.4,2.4);
  \draw[oppgcol!70!black, thick] (2.85,2.4) -- (3.45,2.4);
  % Kanne
  \draw[very thick, black!75] (0,3.0) -- (0,0) -- (2.4,0) -- (2.4,0.3) -- (3.75,3.0);
  \draw[very thick, black!75] (2.4,3.0) -- (2.4,1.5) -- (3.15,3.0);
  % Hank
  \draw[very thick, black!75] (0,2.6) .. controls (-0.9,2.6) and (-0.9,0.8) .. (0,0.8);
  % Punkter
  \draw[dashed, gray!70] (0.6,1.8) -- (2.85,1.8);
  \fill[black] (1.0,1.8) circle (2pt) node[below] {A};
  \fill[black] (2.85,1.8) circle (2pt) node[right=4pt] {B};
\end{tikzpicture}
\end{center}

\begin{enumerate}[a)]
    \item En trykkmåler viser at trykket i A er \SI{1.05e5}{\pascal}. Hva er trykket i B? Begrunn svaret.
    \item Bruk svaret i a) til å forklare hvorfor vannet i tuten står like høyt som vannet i kanna.
    \item Du vipper kanna slik at vannet renner ut av tuten. Kan du fortsatt være sikker på at trykket i A og B er likt? Begrunn svaret.
\end{enumerate}
%
%
}

\svar[title={Løsningsforslag}]{%
\noindent
\textbf{a)} Punktene A og B ligger på samme horisontale nivå, og vi kan gå fra A til B gjennom vannet uten å forlate det. Vannet står i ro. Da er trykket likt i de to punktene (se \textit{Horisontalt væskeelement} i Kapittel~\ref{komp-Chapter:fluiddynamikk}):
\[
    p_B = p_A = \doubleunderline{\SI{1.05e5}{\pascal}}
\]

\smallskip\noindent
\textbf{b)} Trykket i et punkt skyldes vekten av vannet over punktet (pluss lufttrykket). Tenk deg at vannet i tuten stod lavere enn vannet i kanna. Da ville det vært mindre vann over B enn over A, og trykket i B ville vært lavere enn i A. Vannet ville da strømmet fra A mot B, altså inn i tuten, helt til trykkene var like. Vannet kommer derfor bare til ro når det står like høyt i tuten som i kanna.

\smallskip\noindent
\textbf{c)} \doubleunderline{Nei.} Regelen $p_A = p_B$ gjelder bare for en væske i ro. Når vannet strømmer ut gjennom tuten, er væsken i bevegelse, og da kan trykket variere langs et horisontalt nivå. (Det er nettopp en trykkforskjell som får vannet til å strømme.)
}

\oppgaver{Vanntårn}{k7oppg4}{%
%
Mange byer får vannet sitt fra et vanntårn eller et basseng som ligger høyt i terrenget. Vannflaten i et vanntårn ligger \SI{40}{\metre} over en vannkran i et hus. Vannet i ledningen står i ro så lenge kranen er stengt. Lufttrykket er $p_0 = \SI{1.013e5}{\pascal}$.
\begin{enumerate}[a)]
    \item Finn det hydrostatiske trykket ved kranen, altså den delen av trykket som skyldes vannet.
    \item Finn det totale trykket ved kranen.
    \item Et annet hus ligger høyere opp i bakken, slik at kranen der er \SI{25}{\metre} under vannflaten i tårnet. Finn det totale trykket ved denne kranen. Hvilket av husene får best vanntrykk?
\end{enumerate}
%
%
}

\svar[title={Løsningsforslag}]{%
\noindent
\textbf{a)} Det hydrostatiske trykket er $\rho\mathrm{g}h$, der $h$ er dybden under vannflaten (se \textit{Vertikalt væskeelement} i Kapittel~\ref{komp-Chapter:fluiddynamikk}):
\[
    \rho\mathrm{g}h = \SI{1000}{\kilo\gram\per\metre\cubed}\cdot\SI{9.81}{\metre\per\second\squared}\cdot\SI{40}{\metre} = \SI{3.92e5}{\pascal} \approx \doubleunderline{\SI{3.9e5}{\pascal}}
\]

\smallskip\noindent
\textbf{b)} Det totale trykket er lufttrykket pluss det hydrostatiske trykket:
\[
    p = p_0 + \rho\mathrm{g}h = \SI{1.013e5}{\pascal} + \SI{3.92e5}{\pascal} = \SI{4.94e5}{\pascal} \approx \doubleunderline{\SI{4.9e5}{\pascal}}
\]

\smallskip\noindent
\textbf{c)} Nå er dybden under vannflaten \SI{25}{\metre}:
\[
    p = p_0 + \rho\mathrm{g}h = \SI{1.013e5}{\pascal} + \SI{1000}{\kilo\gram\per\metre\cubed}\cdot\SI{9.81}{\metre\per\second\squared}\cdot\SI{25}{\metre}
\]
\[
    p = \SI{1.013e5}{\pascal} + \SI{2.45e5}{\pascal} = \SI{3.47e5}{\pascal} \approx \doubleunderline{\SI{3.5e5}{\pascal}}
\]
Huset som ligger lavest (lengst under vannflaten i tårnet) får det største trykket og dermed best vanntrykk. Det er derfor vanntårn bygges høyt, eller på toppen av en høyde.
}

\oppgaver{Billøfter på verkstedet}{k7oppg5}{%
%
En billøfter på et verksted virker som en hydraulisk jekk. Det lille stempelet har diameter \SI{2.0}{\centi\metre}, og det store stempelet har diameter \SI{20}{\centi\metre}. En bil med masse \SI{1500}{\kilo\gram} står på det store stempelet. Vi ser bort fra massen til stemplene og antar at begge stemplene er i samme høyde.
\begin{enumerate}[a)]
    \item Hvor mange ganger større er arealet av det store stempelet enn arealet av det lille? (Arealet av en sirkel er $A = \pi r^2$.)
    \item Hvor stor kraft må virke på det lille stempelet for å holde bilen oppe?
    \item Det lille stempelet presses \SI{10}{\centi\metre} ned. Hvor mye løftes bilen? \textit{Hint: Væsken som presses ut av den lille sylinderen, må gå inn i den store.}
    \item Regn ut arbeidet som gjøres på det lille stempelet, og arbeidet som gjøres på bilen. Hva ser du? Får vi noe «gratis» med en hydraulisk jekk?
\end{enumerate}
%
%
}

\svar[title={Løsningsforslag}]{%
\noindent
\textbf{a)} Radiene er $r_1 = \SI{1.0}{\centi\metre}$ og $r_2 = \SI{10}{\centi\metre}$. Forholdet mellom arealene er
\[
    \frac{A_2}{A_1} = \frac{\pi r_2^2}{\pi r_1^2} = \left(\frac{r_2}{r_1}\right)^2 = \left(\frac{\SI{10}{\centi\metre}}{\SI{1.0}{\centi\metre}}\right)^2 = \doubleunderline{100}
\]
Når diameteren blir ti ganger så stor, blir arealet hundre ganger så stort.

\smallskip\noindent
\textbf{b)} For å holde bilen oppe må kraften fra væsken på det store stempelet være lik tyngden av bilen, $F_2 = m\mathrm{g} = \SI{1500}{\kilo\gram}\cdot\SI{9.81}{\metre\per\second\squared} = \SI{14715}{\newton}$. Fra Pascals lov (se \textit{Pascals lov} i Kapittel~\ref{komp-Chapter:fluiddynamikk}):
\[
    \frac{F_1}{A_1} = \frac{F_2}{A_2}
    \quad\implies\quad
    F_1 = F_2\cdot\frac{A_1}{A_2} = \frac{\SI{14715}{\newton}}{100} = \SI{147}{\newton} \approx \doubleunderline{\SI{1.5e2}{\newton}}
\]
Det tilsvarer tyngden av omtrent \SI{15}{\kilo\gram}.

\smallskip\noindent
\textbf{c)} Væsken kan ikke presses sammen, så volumet som presses ut av den lille sylinderen, må være lik volumet som går inn i den store:
\[
    A_1 s_1 = A_2 s_2
    \quad\implies\quad
    s_2 = s_1\cdot\frac{A_1}{A_2} = \frac{\SI{10}{\centi\metre}}{100} = \SI{0.10}{\centi\metre} = \doubleunderline{\SI{1.0}{\milli\metre}}
\]

\smallskip\noindent
\textbf{d)} Arbeidet på det lille stempelet:
\[
    W_1 = F_1 s_1 = \SI{147}{\newton}\cdot\SI{0.10}{\metre} = \SI{14.7}{\joule} \approx \doubleunderline{\SI{15}{\joule}}
\]
Arbeidet på bilen:
\[
    W_2 = F_2 s_2 = \SI{14715}{\newton}\cdot\SI{1.0e-3}{\metre} = \SI{14.7}{\joule} \approx \doubleunderline{\SI{15}{\joule}}
\]
Arbeidene er like store. Den hydrauliske jekken gir oss en hundre ganger større kraft, men vi må presse det lille stempelet hundre ganger lenger enn bilen løftes. Vi får altså ikke noe gratis: energien er bevart.
}

\oppgaver{Hydrauliske bremser}{k7oppg6}{%
%
Når du tråkker på bremsepedalen i en bil, presses et lite stempel inn i en sylinder med bremsevæske. Røret med bremsevæske fører videre til et stempel ved hvert av de fire hjulene. Stempelet ved pedalen har areal \SI{1.5}{\centi\metre\squared}, og hvert av stemplene ved hjulene har areal \SI{9.0}{\centi\metre\squared}. Kraften fra pedalen på det lille stempelet er \SI{200}{\newton}.
\begin{enumerate}[a)]
    \item Finn trykkøkningen i bremsevæsken.
    \item Finn kraften fra bremsevæsken på hvert av stemplene ved hjulene.
    \item Rørene fram til hjulene har ulik lengde. Hvorfor blir likevel kraften på stempelet like stor ved alle fire hjulene?
\end{enumerate}
%
%
}

\svar[title={Løsningsforslag}]{%
\noindent
\textbf{a)} Trykkøkningen er kraften delt på arealet til det lille stempelet, $A_1 = \SI{1.5e-4}{\metre\squared}$:
\[
    \Delta p = \frac{F_1}{A_1} = \frac{\SI{200}{\newton}}{\SI{1.5e-4}{\metre\squared}} = \SI{1.33e6}{\pascal} \approx \doubleunderline{\SI{1.3e6}{\pascal}}
\]

\smallskip\noindent
\textbf{b)} Etter Pascals lov får bremsevæsken den samme trykkøkningen overalt. Kraften på hvert stempel ved hjulene, med $A_2 = \SI{9.0e-4}{\metre\squared}$, er da
\[
    F_2 = \Delta p\cdot A_2 = \SI{1.33e6}{\pascal}\cdot\SI{9.0e-4}{\metre\squared} = \doubleunderline{\SI{1.2e3}{\newton}}
\]
Hvert stempel ved hjulene får altså en kraft som er seks ganger så stor som kraften fra pedalen.

\smallskip\noindent
\textbf{c)} Ifølge Pascals lov forplanter trykkøkningen seg uforminsket til alle deler av væsken, uansett hvor lange rørene er og hvilken vei de går. Siden stemplene ved hjulene har samme areal, blir kraften $F = \Delta p\cdot A$ like stor ved alle fire hjulene. Bilen bremser derfor likt på alle hjulene.
}

\oppgaver{Er kronen av gull?}{k7oppg7}{%
%
Ifølge en gammel fortelling fikk Arkimedes i oppdrag av kongen å finne ut om en ny krone var laget av rent gull. Vi gjør et tilsvarende forsøk. En krone har masse \SI{1.00}{\kilo\gram}. Vi henger kronen i en kraftmåler og senker den helt ned i vann. Mens kronen henger i ro under vann, viser kraftmåleren \SI{9.15}{\newton}.
\begin{enumerate}[a)]
    \item Tegn kreftene som virker på kronen mens den henger under vann.
    \item Finn oppdriftskraften på kronen.
    \item Finn volumet til kronen.
    \item Finn massetettheten til kronen. Er den laget av rent gull?
\end{enumerate}
%
%
}

\svar[title={Løsningsforslag}]{%
\noindent
\textbf{a)} Tre krefter virker på kronen: tyngdekraften $\vec{G}$ nedover, oppdriftskraften $\vec{O}$ oppover og kraften $\vec{S}$ fra kraftmåleren oppover. Det er $S$ som kraftmåleren viser.

\smallskip\noindent
\textbf{b)} Kronen henger i ro, så $\sum F = 0$:
\[
    S + O - G = 0 \quad\implies\quad O = G - S
\]
Tyngdekraften er $G = m\mathrm{g} = \SI{1.00}{\kilo\gram}\cdot\SI{9.81}{\metre\per\second\squared} = \SI{9.81}{\newton}$, så
\[
    O = \SI{9.81}{\newton} - \SI{9.15}{\newton} = \doubleunderline{\SI{0.660}{\newton}}
\]

\smallskip\noindent
\textbf{c)} Kronen er helt under vann, så det fortrengte volumet er lik volumet til kronen. Fra $O = \rho_v\mathrm{g}V$ (se \textit{Oppdriftskraft} i Kapittel~\ref{komp-Chapter:fluiddynamikk}):
\[
    V = \frac{O}{\rho_v\mathrm{g}} = \frac{\SI{0.660}{\newton}}{\SI{1000}{\kilo\gram\per\metre\cubed}\cdot\SI{9.81}{\metre\per\second\squared}} = \SI{6.728e-5}{\metre\cubed} \approx \doubleunderline{\SI{6.73e-5}{\metre\cubed}}
\]

\smallskip\noindent
\textbf{d)} Massetettheten er
\[
    \rho = \frac{m}{V} = \frac{\SI{1.00}{\kilo\gram}}{\SI{6.73e-5}{\metre\cubed}} = \SI{1.486e4}{\kilo\gram\per\metre\cubed} \approx \doubleunderline{\SI{1.49e4}{\kilo\gram\per\metre\cubed}}
\]
Rent gull har $\rho = \SI{19.3e3}{\kilo\gram\per\metre\cubed}$. Kronen har lavere tetthet, så den er \doubleunderline{ikke laget av rent gull}. Den er trolig blandet med et lettere metall, for eksempel sølv. En krone av rent gull ville hatt mindre volum, fått mindre oppdrift, og kraftmåleren ville vist omtrent \SI{9.30}{\newton}.
}

\oppgaver{Trekloss som flyter}{k7oppg8}{%
%
En rektangulær trekloss er \SI{20}{\centi\metre} lang, \SI{10}{\centi\metre} bred og \SI{10}{\centi\metre} høy. Treet har massetetthet \SI{600}{\kilo\gram\per\metre\cubed}. Klossen legges i vann, der den flyter med den største flaten ned.
\begin{enumerate}[a)]
    \item Finn massen til klossen.
    \item Hvor stor andel av klossens volum er under vann?
    \item Hvor dypt stikker klossen ned i vannet?
    \item Hvor stor masse kan du legge oppå klossen før den synker helt ned under vannflaten?
\end{enumerate}
%
%
}

\svar[title={Løsningsforslag}]{%
\noindent
\textbf{a)} Volumet til klossen er
\[
    V = \SI{0.20}{\metre}\cdot\SI{0.10}{\metre}\cdot\SI{0.10}{\metre} = \SI{2.0e-3}{\metre\cubed},
\]
og massen er
\[
    m = \rho_L V = \SI{600}{\kilo\gram\per\metre\cubed}\cdot\SI{2.0e-3}{\metre\cubed} = \doubleunderline{\SI{1.2}{\kilo\gram}}
\]

\smallskip\noindent
\textbf{b)} Klossen flyter i ro, så oppdriftskraften er lik tyngdekraften, $O = G$:
\[
    \rho_v\mathrm{g}V_\text{f} = \rho_L\mathrm{g}V
    \quad\implies\quad
    \frac{V_\text{f}}{V} = \frac{\rho_L}{\rho_v} = \frac{\SI{600}{\kilo\gram\per\metre\cubed}}{\SI{1000}{\kilo\gram\per\metre\cubed}} = 0.60 = \doubleunderline{\SI{60}{\percent}}
\]
(Dette er samme fremgangsmåte som i eksempelet \textit{Toppen av isfjellet} i Kapittel~\ref{komp-Chapter:fluiddynamikk}.)

\smallskip\noindent
\textbf{c)} Klossen ligger med den største flaten ned, så høyden er \SI{10}{\centi\metre}. Siden 60\,\% av volumet er under vann, og grunnflaten er den samme hele veien opp, er også 60\,\% av høyden under vann:
\[
    h_\text{f} = 0.60\cdot\SI{10}{\centi\metre} = \doubleunderline{\SI{6.0}{\centi\metre}}
\]

\smallskip\noindent
\textbf{d)} Den største oppdriftskraften klossen kan få, er når hele klossen er under vann, $V_\text{f} = V$:
\[
    O_\text{maks} = \rho_v\mathrm{g}V = \SI{1000}{\kilo\gram\per\metre\cubed}\cdot\SI{9.81}{\metre\per\second\squared}\cdot\SI{2.0e-3}{\metre\cubed} = \SI{19.6}{\newton}
\]
Denne kraften kan bære en total masse på
\[
    m_\text{total} = \frac{O_\text{maks}}{\mathrm{g}} = \frac{\SI{19.6}{\newton}}{\SI{9.81}{\metre\per\second\squared}} = \SI{2.0}{\kilo\gram}
\]
Klossen selv har masse \SI{1.2}{\kilo\gram}, så vi kan legge på
\[
    \Delta m = \SI{2.0}{\kilo\gram} - \SI{1.2}{\kilo\gram} = \doubleunderline{\SI{0.80}{\kilo\gram}}
\]
}

\oppgaver{Lufta i et klasserom}{k7oppg9}{%
%
Et klasserom er \SI{5.0}{\metre} langt, \SI{4.0}{\metre} bredt og \SI{2.5}{\metre} høyt. Lufta har massetetthet \SI{1.2}{\kilo\gram\per\metre\cubed}.
\begin{enumerate}[a)]
    \item Finn massen til lufta i klasserommet.
    \item Sammenlign med massen til en voksen person. Er svaret overraskende?
\end{enumerate}
%
%
}

\svar[title={Løsningsforslag}]{%
\noindent
\textbf{a)} Volumet av rommet er
\[
    V = \SI{5.0}{\metre}\cdot\SI{4.0}{\metre}\cdot\SI{2.5}{\metre} = \SI{50}{\metre\cubed},
\]
og massen til lufta er
\[
    m = \rho V = \SI{1.2}{\kilo\gram\per\metre\cubed}\cdot\SI{50}{\metre\cubed} = \doubleunderline{\SI{60}{\kilo\gram}}
\]

\smallskip\noindent
\textbf{b)} Lufta i klasserommet har omtrent like stor masse som en voksen person. Mange blir overrasket over dette, fordi vi ikke merker at lufta veier noe. Lufta har liten massetetthet, men et klasserom har et stort volum.
}

\oppgaver{Dykking i sjøen}{k7oppg10}{%
%
En dykker svømmer \SI{30}{\metre} under overflaten i sjøen. Sjøvann har massetetthet \SI{1030}{\kilo\gram\per\metre\cubed}, og lufttrykket ved overflaten er $p_0 = \SI{1.013e5}{\pascal}$.
\begin{enumerate}[a)]
    \item Finn det totale trykket på dykkeren.
    \item Hvor mange ganger større er dette trykket enn lufttrykket ved overflaten?
\end{enumerate}
%
%
}

\svar[title={Løsningsforslag}]{%
\noindent
\textbf{a)} Vi bruker $p = p_0 + \rho\mathrm{g}h$:
\[
    p = \SI{1.013e5}{\pascal} + \SI{1030}{\kilo\gram\per\metre\cubed}\cdot\SI{9.81}{\metre\per\second\squared}\cdot\SI{30}{\metre}
\]
\[
    p = \SI{1.013e5}{\pascal} + \SI{3.03e5}{\pascal} = \SI{4.04e5}{\pascal} \approx \doubleunderline{\SI{4.0e5}{\pascal}}
\]

\smallskip\noindent
\textbf{b)}
\[
    \frac{p}{p_0} = \frac{\SI{4.04e5}{\pascal}}{\SI{1.013e5}{\pascal}} = 3.99 \approx \doubleunderline{4.0}
\]
Trykket på \SI{30}{\metre} dybde er omtrent fire ganger så stort som ved overflaten: én atmosfære fra lufta og omtrent én atmosfære for hver tiende meter med vann.
}

\oppgaver{Sugerør}{k7oppg11}{%
%
Når du drikker med sugerør, senker du trykket i munnen. Da er det lufttrykket på væskeoverflaten i glasset som presser væsken opp i sugerøret.
\begin{enumerate}[a)]
    \item Tenk deg at du klarer å fjerne all lufta i sugerøret, slik at trykket over vannet i røret blir null. Hvor høyt kan lufttrykket da presse vannet opp i røret? Lufttrykket er $p_0 = \SI{1.013e5}{\pascal}$.
    \item Kan man drikke vann med et \SI{12}{\metre} langt sugerør, uansett hvor kraftig man suger? Begrunn svaret.
\end{enumerate}
%
%
}

\svar[title={Løsningsforslag}]{%
\noindent
\textbf{a)} Vi ser på vannsøylen i sugerøret. Ved vannflaten i glasset er trykket $p_0$. Inne i røret, i samme høyde, er trykket også $p_0$ (samme nivå i den samme sammenhengende væsken). Over vannet i røret er trykket null. Trykkdifferansen mellom bunnen og toppen av vannsøylen i røret er da
\[
    p_0 - 0 = \rho\mathrm{g}h
    \quad\implies\quad
    h = \frac{p_0}{\rho\mathrm{g}} = \frac{\SI{1.013e5}{\pascal}}{\SI{1000}{\kilo\gram\per\metre\cubed}\cdot\SI{9.81}{\metre\per\second\squared}} = \doubleunderline{\SI{10.3}{\metre}}
\]

\smallskip\noindent
\textbf{b)} \doubleunderline{Nei.} Selv med et perfekt vakuum i røret kan lufttrykket bare presse vannet \SI{10.3}{\metre} opp. Det hjelper ikke å suge hardere, for vi kan ikke senke trykket under null. Det er lufttrykket som gjør jobben, ikke selve sugingen.
}

\oppgaver{Vindu i en ubåt}{k7oppg12}{%
%
En ubåt ligger \SI{100}{\metre} under overflaten i sjøen ($\rho = \SI{1030}{\kilo\gram\per\metre\cubed}$). Inne i ubåten er trykket det samme som lufttrykket ved overflaten. Ubåten har et vindu med areal \SI{0.050}{\metre\squared}.
\begin{enumerate}[a)]
    \item Finn trykkforskjellen mellom utsiden og innsiden av vinduet.
    \item Finn nettokraften fra vannet og lufta på vinduet.
\end{enumerate}
\textit{Hint: Lufttrykket virker både på overflaten av sjøen og inne i ubåten.}
%
%
}

\svar[title={Løsningsforslag}]{%
\noindent
\textbf{a)} Utenfor vinduet er trykket $p_\text{ute} = p_0 + \rho\mathrm{g}h$, og inne i ubåten er det $p_\text{inne} = p_0$. Trykkforskjellen er derfor bare det hydrostatiske trykket:
\[
    \Delta p = p_\text{ute} - p_\text{inne} = \rho\mathrm{g}h = \SI{1030}{\kilo\gram\per\metre\cubed}\cdot\SI{9.81}{\metre\per\second\squared}\cdot\SI{100}{\metre} = \SI{1.01e6}{\pascal} \approx \doubleunderline{\SI{1.0e6}{\pascal}}
\]

\smallskip\noindent
\textbf{b)} Nettokraften på vinduet er trykkforskjellen ganget med arealet:
\[
    F = \Delta p\cdot A = \SI{1.01e6}{\pascal}\cdot\SI{0.050}{\metre\squared} = \SI{5.05e4}{\newton} \approx \doubleunderline{\SI{5.1e4}{\newton}}
\]
Det tilsvarer tyngden av omtrent fem tonn. Vinduet må derfor være svært kraftig.
}

\oppgaver{Heliumballong}{k7oppg13}{%
%
Arkimedes' lov gjelder ikke bare i væsker, men også i gasser. En ballong er fylt med \SI{0.010}{\metre\cubed} helium. Helium har massetetthet \SI{0.18}{\kilo\gram\per\metre\cubed}, og lufta rundt har massetetthet \SI{1.2}{\kilo\gram\per\metre\cubed}. Selve ballongen (gummien og snora) har masse \SI{3.0}{\gram}. Vi ser bort fra volumet til gummien.
\begin{enumerate}[a)]
    \item Finn oppdriftskraften fra lufta på ballongen.
    \item Finn tyngdekraften på ballongen med helium.
    \item Vil ballongen stige eller synke? Finn nettokraften på ballongen.
\end{enumerate}
%
%
}

\svar[title={Løsningsforslag}]{%
\noindent
\textbf{a)} Ballongen fortrenger \SI{0.010}{\metre\cubed} luft. Etter Arkimedes' lov er oppdriftskraften lik tyngden av denne lufta:
\[
    O = \rho_\text{luft}\mathrm{g}V = \SI{1.2}{\kilo\gram\per\metre\cubed}\cdot\SI{9.81}{\metre\per\second\squared}\cdot\SI{0.010}{\metre\cubed} = \SI{0.118}{\newton} \approx \doubleunderline{\SI{0.12}{\newton}}
\]

\smallskip\noindent
\textbf{b)} Massen til heliumet er $m_\text{He} = \rho_\text{He}V = \SI{0.18}{\kilo\gram\per\metre\cubed}\cdot\SI{0.010}{\metre\cubed} = \SI{1.8e-3}{\kilo\gram}$. Den totale massen er
\[
    m = \SI{1.8e-3}{\kilo\gram} + \SI{3.0e-3}{\kilo\gram} = \SI{4.8e-3}{\kilo\gram},
\]
og tyngdekraften er
\[
    G = m\mathrm{g} = \SI{4.8e-3}{\kilo\gram}\cdot\SI{9.81}{\metre\per\second\squared} = \SI{0.0471}{\newton} \approx \doubleunderline{\SI{0.047}{\newton}}
\]

\smallskip\noindent
\textbf{c)} Med positiv retning oppover er nettokraften
\[
    F_\text{netto} = O - G = \SI{0.118}{\newton} - \SI{0.047}{\newton} = \SI{0.071}{\newton}
\]
Nettokraften er positiv, så ballongen \doubleunderline{stiger}, med en nettokraft på omtrent \doubleunderline{\SI{0.071}{\newton}} oppover.
}

\oppgaver{Olje og vann i et U-rør}{k7oppg14}{%
%
Et U-rør inneholder vann. I det venstre røret heller vi oppå vannet en \SI{10.0}{\centi\metre} høy søyle med olje, som har massetetthet \SI{800}{\kilo\gram\per\metre\cubed}. Olje og vann blander seg ikke. Når alt står i ro, ligger grenseflaten mellom olje og vann i den venstre delen av røret (se figuren).

\begin{center}
\begin{tikzpicture}[scale=0.9, font=\small]
  % Vann
  \fill[oppgcol!25] (0,0) -- (2.6,0) -- (2.6,1.8) -- (2.0,1.8) -- (2.0,0.6) -- (0.6,0.6) -- (0.6,1.0) -- (0,1.0) -- cycle;
  % Olje
  \fill[yellow!45!orange!40] (0,1.0) rectangle (0.6,2.0);
  % Rør
  \draw[very thick, black!75] (0,3.0) -- (0,0) -- (2.6,0) -- (2.6,3.0);
  \draw[very thick, black!75] (0.6,3.0) -- (0.6,0.6) -- (2.0,0.6) -- (2.0,3.0);
  % Nivålinje ved grenseflaten
  \draw[dashed, gray!60!black] (-0.3,1.0) -- (2.9,1.0);
  \fill[black] (0.3,1.0) circle (1.8pt) node[below right=-1pt, font=\footnotesize] {A};
  \fill[black] (2.3,1.0) circle (1.8pt) node[below right=-1pt, font=\footnotesize] {B};
  % Etiketter
  \draw[thin, black!60] (0.45,1.7) -- (1.2,2.3) node[right, font=\footnotesize] {olje};
  \node[right, font=\footnotesize] at (2.7,1.4) {vann};
  \draw[<->, thin, gray!40!black] (-0.3,1.0) -- (-0.3,2.0)
      node[midway, left, font=\footnotesize] {\SI{10.0}{\centi\metre}};
\end{tikzpicture}
\end{center}

\begin{enumerate}[a)]
    \item Punkt A ligger ved grenseflaten mellom olje og vann i det venstre røret, og punkt B ligger i samme høyde i det høyre røret. Forklar hvorfor trykket i A og B er likt.
    \item Hvor høyt over grenseflaten står vannet i det høyre røret?
    \item Hvor mye høyere står oljeoverflaten enn vannoverflaten?
\end{enumerate}
%
%
}

\svar[title={Løsningsforslag}]{%
\noindent
\textbf{a)} Under punkt A er det bare vann, og vannet henger sammen gjennom bunnen av U-røret og opp til punkt B. A og B ligger altså på samme horisontale nivå i den samme sammenhengende væsken (vannet), og alt står i ro. Da er $p_A = p_B$.

(Over grenseflaten er det olje i det venstre røret og vann i det høyre. Der er det ikke samme væske på begge sider, og da kan vi ikke bruke denne regelen.)

\smallskip\noindent
\textbf{b)} Trykket i A skyldes lufta og oljesøylen over A. Trykket i B skyldes lufta og vannsøylen over B:
\[
    p_A = p_0 + \rho_\text{olje}\mathrm{g}h_\text{olje}
    \qquad\text{og}\qquad
    p_B = p_0 + \rho_\text{vann}\mathrm{g}h_\text{vann}
\]
Vi setter $p_A = p_B$. Da faller $p_0$ og $\mathrm{g}$ bort, og vi får
\[
    \rho_\text{olje}h_\text{olje} = \rho_\text{vann}h_\text{vann}
    \quad\implies\quad
    h_\text{vann} = \frac{\rho_\text{olje}}{\rho_\text{vann}}\,h_\text{olje} = \frac{800}{1000}\cdot\SI{10.0}{\centi\metre} = \doubleunderline{\SI{8.0}{\centi\metre}}
\]

\smallskip\noindent
\textbf{c)} Oljeoverflaten ligger \SI{10.0}{\centi\metre} over grenseflaten, og vannoverflaten ligger \SI{8.0}{\centi\metre} over grenseflaten. Forskjellen er
\[
    \SI{10.0}{\centi\metre} - \SI{8.0}{\centi\metre} = \doubleunderline{\SI{2.0}{\centi\metre}}
\]
Oljen står høyest fordi den har lavere massetetthet enn vann. Det trengs en høyere oljesøyle for å gi det samme trykket.
}
